\begin{salyga}
Išspręskite lygtį:
    \[ (x - 1)(x + 2)(x + 3)(x + 6) = 160 \]
\end{salyga}

\begin{atsakymas}
$x=-7$ arba $x=2$.
\end{atsakymas}

\begin{sprendimas}
Sugrupuokime pirmąjį su ketvirtuoju, o antrąjį su trečiuoju daugikliu:
\[
(x-1)(x+6)=x^2+5x-6,\qquad
(x+2)(x+3)=x^2+5x+6.
\]
Tegu $t=x^2+5x$. Tada $(t-6)(t+6)=160$, todėl $t^2=196$ ir $t=14$ arba $t=-14$.

Kai $t=14$, $(x+7)(x-2)=0$, taigi $x=-7$ arba $x=2$. Kai $t=-14$, lygties $x^2+5x+14=0$ diskriminantas yra $25-56=-31<0$, todėl realiųjų sprendinių nėra.
\end{sprendimas}
